\documentclass[UTF8,12pt,a4paper,fontset=fandol]{ctexart}

\usepackage[a4paper,top=25mm,bottom=25mm,left=27mm,right=27mm]{geometry}
\usepackage{amsmath,amssymb,bm}
\usepackage{enumitem}
\usepackage{fancyhdr}
\usepackage{microtype}

\setlength{\parindent}{2em}
\setlength{\parskip}{0.25em}
\setlist[enumerate]{leftmargin=2.2em,itemsep=0.25em,topsep=0.35em}
\allowdisplaybreaks[2]

\pagestyle{fancy}
\setlength{\headheight}{15pt}
\fancyhf{}
\fancyhead[C]{二维扩散方程数值测试}
\fancyfoot[C]{\thepage}
\renewcommand{\headrulewidth}{0.4pt}
\renewcommand{\footrulewidth}{0pt}
\makeatletter
\let\ps@plain\ps@fancy
\makeatother

\ctexset{
  section = {
    format = \Large\bfseries,
    number = \arabic{section},
    aftername = \quad
  },
  subsection = {
    format = \large\bfseries,
    number = \thesection.\arabic{subsection},
    aftername = \quad
  },
  subsubsection = {
    format = \normalsize\bfseries,
    number =,
    aftername =
  }
}

\begin{document}

\begin{center}
  {\LARGE\bfseries 第二题：二维扩散方程}
\end{center}
\vspace{0.8em}

\section{扩散方程}

\subsection{控制方程}

考虑二维扩散方程

\begin{equation}
\frac{\partial\phi}{\partial t}
-\nabla\cdot\left(\mu\nabla\phi\right)=0,
\end{equation}

其中，$\phi=\phi(x,y,t)$ 为待求标量，$\mu$ 为扩散系数。

请基于有限体积法写出该方程的半离散形式，并分别采用显式时间离散格式和显式空间离散格式构建数值求解器。

\subsection{数值算例}

\subsubsection{2.2.1\quad 具有间断初值的扩散问题}

考虑扩散系数

\begin{equation}
\mu=1
\end{equation}

的扩散方程。其精确解为

\begin{align}
\phi(x,y,t)
=\frac14
&\left[
-\operatorname{Erf}\left(\frac{-1-x}{2\sqrt{t}}\right)
+\operatorname{Erf}\left(\frac{1-x}{2\sqrt{t}}\right)
\right]
\notag\\
&\times
\left[
-\operatorname{Erf}\left(\frac{-1-y}{2\sqrt{t}}\right)
+\operatorname{Erf}\left(\frac{1-y}{2\sqrt{t}}\right)
\right],
\end{align}

其中误差函数定义为

\begin{equation}
\operatorname{Erf}(z)
=\frac{2}{\sqrt{\pi}}
\int_0^z e^{-s^2}\,\mathrm ds.
\end{equation}

上述精确解在 $t\to0^+$ 时对应的间断初始条件为

\begin{equation}
\phi(x,y,0)=
\begin{cases}
1, & |x|\leq1,\quad |y|\leq1,\\[1mm]
0, & \text{其他位置}.
\end{cases}
\end{equation}

计算区域为

\begin{equation}
\Omega=[-5,5]^2.
\end{equation}

在区域边界上施加齐次 Neumann 边界条件

\begin{equation}
\frac{\partial\phi}{\partial n}=0,
\qquad (x,y)\in\partial\Omega.
\end{equation}

分别采用三角形单元和四边形单元对计算区域进行划分，并逐渐加密网格。

请在

\begin{equation}
t=0.2
\end{equation}

时计算数值误差和收敛阶，以评价数值解的精度。

\subsubsection{2.2.2\quad 连续函数扩散问题}

考虑扩散系数

\begin{equation}
\mu=1
\end{equation}

的二维扩散方程。其连续精确解为

\begin{equation}
\phi(x,y,t)
=\frac{1}{1+4t}
\exp\left[
-\frac{\mu(x^2+y^2)}{1+4t}
\right].
\end{equation}

对应的初始条件为

\begin{equation}
\phi(x,y,0)
=\exp\left[-\mu(x^2+y^2)\right]
=\exp\left[-(x^2+y^2)\right].
\end{equation}

计算区域为

\begin{equation}
\Omega=[-5,5]^2.
\end{equation}

在区域边界上施加由精确解给出的 Dirichlet 边界条件

\begin{equation}
\phi(x,y,t)
=\frac{1}{1+4t}
\exp\left[
-\frac{x^2+y^2}{1+4t}
\right],
\qquad (x,y)\in\partial\Omega.
\end{equation}

分别采用三角形单元和四边形单元对计算区域进行划分，并逐渐加密网格。所有计算中的时间步长均根据显式扩散格式的稳定性条件进行调整。

请在

\begin{equation}
t=0.2
\end{equation}

时分析数值结果的精度和收敛阶。

\end{document}
